% \section{Counterfactual Model Probing}
% \label{sec:probing}

% The fidelity results enable a compact downstream protocol. This section specifies the endpoints without inserting unmeasured projected values into the reviewer manuscript. Tables~\ref{tab:model_interpretability} and~\ref{tab:aopc} reserve the paired model-sensitivity and explanation-faithfulness results.

% \begin{table*}[t]
% \centering
% \caption{Model performance and interpretability intervention results across tasks.
% Intervention values are reported as TVD / Flip Rate.}
% \label{tab:xai_all_tasks}

% \small
% \renewcommand{\arraystretch}{1.30}
% \setlength{\tabcolsep}{4.5pt}

% \begin{adjustbox}{max width=\textwidth}
% \begin{tabular}{lccccccc}
% \toprule

% \textbf{Model} &
% \textbf{Performance} &
% \shortstack{\textbf{Stain}\\\textbf{Brightness}} &
% \shortstack{\textbf{Nuclear}\\\textbf{Enlargement}} &
% \shortstack{\textbf{Nuclear Shape}\\\textbf{Irregularity}} &
% \shortstack{\textbf{Immune}\\\textbf{Burden}} &
% \shortstack{\textbf{Tumor-Immune}\\\textbf{Mixing}} &
% \textbf{BNR} \\

% \midrule

% \multicolumn{8}{l}{\textbf{PAM50 Classification} \hfill \textit{Performance metric: AUC}} \\
% \midrule

% ResNet-50
% & 0.7627
% & 0.0243 / 0.0238
% & 0.0257 / 0.0281
% & 0.0270 / 0.0241
% & 0.0442 / 0.0626
% & 0.0275 / 0.0267
% & 1.2776 \\

% CTransPath
% & 0.8143
% & 0.0104 / 0.0158
% & 0.0093 / 0.0081
% & 0.0094 / 0.0030
% & 0.0152 / 0.0189
% & 0.0105 / 0.0075
% & 1.0694 \\

% UNI2-h
% & 0.8713
% & 0.0060 / 0.0094
% & 0.0066 / 0.0115
% & 0.0065 / 0.0073
% & 0.0107 / 0.0137
% & 0.0065 / 0.0065
% & 1.2680 \\

% \midrule

% \multicolumn{8}{l}{\textbf{Overall Survival} \hfill \textit{Performance metric: C-index}} \\
% \midrule

% ResNet-50
% & 0.5445
% & 0.0026 / 0.0024
% & 0.0025 / 0.0017
% & 0.0033 / 0.0003
% & 0.0039 / 0.0000
% & 0.0029 / 0.0005
% & 1.1893 \\

% CTransPath
% & 0.6216
% & 0.0017 / 0.0014
% & 0.0017 / 0.0017
% & 0.0018 / 0.0000
% & 0.0026 / 0.0005
% & 0.0015 / 0.0010
% & 1.1398 \\

% UNI2-h
% & 0.6209
% & 0.0014 / 0.0021
% & 0.0015 / 0.0021
% & 0.0014 / 0.0010
% & 0.0023 / 0.0021
% & 0.0014 / 0.0003
% & 1.1711 \\

% PathLUPI + CONCH
% & 0.5860
% & 0.0043 / 0.0057
% & 0.0049 / 0.0032
% & 0.0043 / 0.0043
% & 0.0207 / 0.0207
% & 0.0036 / 0.0043
% & 1.9488 \\

% \bottomrule
% \end{tabular}
% \end{adjustbox}

% \end{table*}

% \subsection{Stain invariance}

% Red, green, and blue controls are swept within their supported ranges while spatial structure, nuclear geometry, and noise remain fixed. RGB interventions are complemented by optical-density measurements because raw channels mix stain concentration with acquisition and tissue composition. For a representation encoder, sensitivity is the normalized embedding displacement per recovered color change. For a task model, it is the change in logit, probability, or risk. A robust model need not be completely color-invariant---hematoxylin changes real nuclear contrast---but nuisance response should be reported relative to verified morphology and spatial response.

% \subsection{Morphology sensitivity}

% Nuclear size, eccentricity, and solidity are swept independently. An intervention enters analysis only when the recovered feature is ordered correctly and non-target controls remain within frozen tolerances. Across-image correlations answer whether natural condition variation is preserved; same-image dose curves answer whether the model responds to a targeted change. We report standardized response, monotonicity, and a signal-to-seed-variation ratio. Using the same cases, doses, and seeds for each model makes the comparison paired.

% \subsection{Spatial sensitivity}

% Spatial maps vary immune density, tumor--immune distance, mixing, clustering, and exclusion while holding morphology and stain fixed. Patch encoders are evaluated by representation movement and linear recovery of intervention magnitude. End-to-end breast models are evaluated at their native output: subtype probability for patch classifiers, slide score for multiple-instance models, and risk for survival systems. Patch interventions are propagated through the same sampling and aggregation process used by the downstream model.

% \subsection{Explanation faithfulness}

% Given a saliency or attention method, regions or cellular concepts are ranked for each source. \method{} applies a matched intervention to top-ranked, random, and bottom-ranked concepts under the same area and magnitude budget. Faithfulness is supported when top-ranked edits produce larger target-output changes than random and bottom-ranked edits, after conditioning on recovered fidelity. This is a model-causal test under the generator; it is not evidence of biological or treatment causality.

% The evaluated panel should distinguish general-purpose encoders (ResNet50, CTransPath, CONCH, UNI2-h, and Virchow) from task-trained breast pathology models. Reporting embedding movement as though it were a clinical prediction would otherwise conflate representation sensitivity with downstream decision behavior.
